% TODO checklist:
% - [x] Inspected src/security/pii_scrubber.py (PII detection and redaction)
% - [x] Inspected src/security/output_guard.py (Cypher safety validation)
% - [x] Inspected src/security/intent_filter.py (Dangerous operation blocking)
% - [x] Reviewed Q&A.md (Data protection patterns)
% - [x] Reviewed README.md (Privacy features overview)

\section{Data Protection}
\label{sec:data-protection}

The Cineca Agentic Platform implements multiple layers of data protection to safeguard sensitive information throughout the request lifecycle. These protections operate at the input, processing, and output stages to prevent data leakage, unauthorized operations, and exposure of personally identifiable information (PII).

\subsection{PII Scrubbing}
\label{subsec:pii-scrubbing}

The PII scrubber module provides zero-dependency heuristics to detect and sanitize personally identifiable information in free text and structured payloads.

\subsubsection{Detection Patterns}

The system detects common PII categories using regular expression patterns:

\begin{lstlisting}[language=Python,caption={PII detection patterns}]
# Email addresses
EMAIL_RE = re.compile(
    r"\b[A-Za-z0-9._%+\-]+@[A-Za-z0-9.\-]+\.[A-Za-z]{2,}\b"
)

# Phone numbers (international format support)
PHONE_RE = re.compile(r"""(?x)
    (?:(?:(?:\+|00)\d{1,3}[\s.\-]?)?  # country code
       (?:\(?\d{2,4}\)?[\s.\-]?)?     # area code
       \d{3,4}[\s.\-]?\d{3,4}(?!\d))  # local number
""")

# IPv4 addresses
IPV4_RE = re.compile(r"\b(?:(?:\d{1,3}\.){3}\d{1,3})\b")

# US Social Security Numbers
SSN_US_RE = re.compile(r"\b\d{3}-\d{2}-\d{4}\b")

# International Bank Account Numbers
IBAN_RE = re.compile(r"\b[A-Z]{2}\d{2}[A-Z0-9]{11,30}\b")

# Credit card numbers (13-19 digits, Luhn validated)
CC_RAW_RE = re.compile(r"\b(?:\d[ \-]?){13,19}\b")
\end{lstlisting}

\subsubsection{Credit Card Validation}

Credit card detection includes Luhn checksum validation to reduce false positives:

\begin{lstlisting}[language=Python,caption={Luhn algorithm for credit card validation}]
def _luhn_valid(digits: str) -> bool:
    """Return True if numeric string passes Luhn check."""
    nums = [int(c) for c in digits if c.isdigit()]
    if len(nums) < 13 or len(nums) > 19:
        return False
    
    checksum = 0
    parity = (len(nums) - 2) % 2
    for i, n in enumerate(nums[:-1]):
        if i % 2 == parity:
            n *= 2
            if n > 9:
                n -= 9
        checksum += n
    
    return (checksum + nums[-1]) % 10 == 0
\end{lstlisting}

\subsubsection{Redaction Modes}

The scrubber supports multiple redaction strategies configured via \texttt{PII\_SCRUBBER\_MODE}:

\begin{table}[ht]
\centering
\caption{PII Scrubber Redaction Modes}
\label{tab:pii-modes}
\begin{tabular}{lp{8cm}}
\hline
\textbf{Mode} & \textbf{Behavior} \\
\hline
\texttt{mask} & Replace with partially masked values (e.g., \texttt{[REDACTED]}) \\
\texttt{hash} & Replace with stable SHA-256 digest (\texttt{sha256:<hex>}) \\
\texttt{remove} & Replace values with \texttt{None} or empty string \\
\texttt{off} & No-op; return input unchanged \\
\hline
\end{tabular}
\end{table}

\begin{lstlisting}[language=Python,caption={PII redaction with mode selection}]
def _replacement(hit: PiiHit, mode: str) -> str:
    if mode == "hash":
        return f"sha256:{_sha256_hex(hit.value)}"
    if mode == "remove":
        return ""
    # mask mode
    if hit.category == "email":
        return "[REDACTED]"
    if hit.category == "phone":
        return "[REDACTED]"
    if hit.category == "credit_card":
        # Keep first/last 4 digits for reference
        digits = re.sub(r"\D", "", hit.value)
        return f"{digits[:4]}{'*'*(len(digits)-8)}{digits[-4:]}"
    return "[REDACTED]"
\end{lstlisting}

\subsubsection{Sensitive Key Detection}

Beyond pattern matching, the scrubber identifies sensitive data by key names:

\begin{lstlisting}[language=Python,caption={Sensitive key detection}]
SENSITIVE_KEYS = {
    "password", "passwd", "secret", "api_key", "apikey",
    "token", "access_token", "refresh_token",
    "authorization", "auth", "ssn", "iban",
    "credit_card", "card_number", "cvv",
    "email", "phone", "tax_id", "national_id",
    "passport", "address", "dob", "birthdate"
}

def scrub_dict(d: Mapping[str, Any], mode: str) -> dict:
    """Scrub a mapping; sensitive keys are redacted."""
    out = {}
    for k, v in d.items():
        if k.lower() in SENSITIVE_KEYS:
            out[k] = _redact_by_keys(k, v, mode)
        else:
            out[k] = scrub(v, mode=mode)  # Recursive
    return out
\end{lstlisting}

\subsubsection{API Surface}

The PII scrubber provides a clean public API:

\begin{lstlisting}[language=Python,caption={PII scrubber public API}]
# Text scrubbing
scrub_text(text: str, mode: str = None) -> str

# Recursive object scrubbing (dict/list/tuple/str)
scrub(obj: Any, mode: str = None) -> Any

# Dictionary-specific scrubbing with key sensitivity
scrub_dict(d: Mapping, mode: str = None) -> dict

# Detection without modification
find_pii(text: str) -> list[dict]  # Returns offsets and categories
contains_pii(text: str) -> bool
\end{lstlisting}

\subsection{Output Guard}
\label{subsec:output-guard}

The output guard module provides safety controls for LLM-generated content, with a particular focus on Cypher query validation.

\subsubsection{Cypher Analysis}

The guard analyzes Cypher queries to detect potentially dangerous operations:

\begin{lstlisting}[language=Python,caption={Cypher safety analysis}]
@dataclass(frozen=True)
class CypherAnalysis:
    text: str
    has_return: bool
    has_limit: bool
    writes: bool           # CREATE, MERGE, SET, DELETE, etc.
    destructive: bool      # DROP GRAPH, TRUNCATE
    unbounded: bool        # Variable-length traversals -[*]->
    risky_call: bool       # CALL with write operations
    risk_score: int        # 0-100 composite score
    reasons: list[str]

def analyze_cypher(query: str) -> CypherAnalysis:
    reasons = []
    risk = 0
    
    writes = bool(RE_WRITE.search(query))
    if writes:
        reasons.append("write-verb present")
        risk += 35
    
    destructive = bool(RE_DROP_GRAPH.search(query))
    if destructive:
        reasons.append("destructive verb present")
        risk += 50
    
    unbounded = bool(RE_UNBOUNDED.search(query))
    if unbounded:
        reasons.append("unbounded variable-length traversal")
        risk += 20
    
    if has_return and not has_limit:
        reasons.append("RETURN without LIMIT")
        risk += 15
    
    return CypherAnalysis(
        text=query,
        risk_score=min(100, risk),
        reasons=reasons,
        # ... other fields
    )
\end{lstlisting}

\subsubsection{Guard Modes}

The output guard operates in three modes:

\begin{table}[ht]
\centering
\caption{Output Guard Operating Modes}
\label{tab:output-guard-modes}
\begin{tabular}{lp{9cm}}
\hline
\textbf{Mode} & \textbf{Behavior} \\
\hline
\texttt{enforce} & Block dangerous queries, auto-append LIMIT, raise HTTP 400 on violations \\
\texttt{monitor} & Log and annotate risks but never block; callers decide \\
\texttt{off} & No safety checks applied \\
\hline
\end{tabular}
\end{table}

\subsubsection{Automatic Limit Injection}

To prevent unbounded result sets, the guard can automatically append LIMIT clauses:

\begin{lstlisting}[language=Python,caption={Automatic LIMIT injection}]
def ensure_cypher_limit(query: str, limit: int = 100) -> str:
    """Append LIMIT N if the query returns rows but has no LIMIT."""
    if not re.search(r"RETURN\b(?!.*\bLIMIT\b)", query, re.I | re.S):
        return query
    
    q_no_semi = re.sub(r"\s*;+$", "", query)
    return f"{q_no_semi} LIMIT {limit}"
\end{lstlisting}

\subsubsection{Guard Configuration}

\begin{table}[ht]
\centering
\caption{Output Guard Configuration Settings}
\label{tab:output-guard-config}
\begin{tabular}{lll}
\hline
\textbf{Setting} & \textbf{Default} & \textbf{Description} \\
\hline
\texttt{OUTPUT\_GUARD\_MODE} & \texttt{monitor} & Operating mode \\
\texttt{OUTPUT\_GUARD\_ALLOW\_WRITES} & \texttt{False} & Allow write operations \\
\texttt{OUTPUT\_GUARD\_ENFORCE\_LIMIT} & \texttt{True} & Auto-append LIMIT \\
\texttt{OUTPUT\_GUARD\_DEFAULT\_LIMIT} & \texttt{100} & Default LIMIT value \\
\texttt{OUTPUT\_GUARD\_BLOCK\_DROP\_GRAPH} & \texttt{True} & Block DROP GRAPH \\
\hline
\end{tabular}
\end{table}

\subsection{Intent Filter}
\label{subsec:intent-filter}

The intent filter provides lightweight guardrails against unsafe or costly requests by analyzing user input for high-risk patterns.

\subsubsection{Risk Categories}

The filter detects multiple threat categories:

\begin{lstlisting}[language=Python,caption={Intent filter risk categories}]
# Prompt injection / jailbreak attempts
RE_PROMPT_INJECTION = re.compile(
    r"(?i)\b(ignore|bypass|override|forget)\b.*"
    r"\b(instruction|policy|guard|safety)\b"
)

# Secrets scraping attempts
RE_SECRETS = re.compile(
    r"(?i)\b(api[_ -]?key|secret|password|token|ssh[_ -]?key)\b"
)

# PII hunting
RE_PII = re.compile(
    r"(?i)\b(ssn|social security|credit\s*card|cvv|iban|"
    r"passport|national id|tax id)\b"
)

# Dangerous shell/OS operations
RE_SHELL = re.compile(
    r"(?i)\b(rm\s+-rf\b|mkfs\w*\b|dd\s+if=\b|"
    r"format\s+c:\b|shutdown\b|poweroff\b|reboot\b)\b"
)

# SQL destructive operations
RE_SQL_DROP = re.compile(
    r"(?i)\b(drop\s+(database|table)\b|truncate\s+table\b)"
)

# Cypher destructive operations
RE_CYPHER_DANGER = re.compile(
    r"(?i)\b(detach\s+delete\b|drop\s+graph\b)"
)

# Bulk exfiltration attempts
RE_EXFIL = re.compile(
    r"(?i)\b(dump|export|download)\b.*"
    r"\b(all|everything|entire|database|db)\b"
)
\end{lstlisting}

\subsubsection{Risk Scoring}

Each detected pattern contributes to a cumulative risk score:

\begin{lstlisting}[language=Python,caption={Intent filter risk scoring}]
def analyze_intent(text: str) -> IntentResult:
    risk = 0
    categories = []
    reasons = []
    
    def add(cat: str, reason: str, weight: int):
        nonlocal risk
        categories.append(cat)
        reasons.append(reason)
        risk = min(100, risk + weight)
    
    if RE_PROMPT_INJECTION.search(text):
        add("prompt_injection", 
            "attempt to override safety", 20)
    
    if RE_SECRETS.search(text):
        add("secrets", 
            "mentions secrets/tokens/passwords", 25)
    
    if RE_SHELL.search(text):
        add("system_abuse", 
            "dangerous shell operation", 40)
    
    if RE_SQL_DROP.search(text):
        # Immediate block for destructive SQL
        return IntentResult(
            allowed=False,
            action="block",
            risk_score=100,
            categories=["db_destructive"],
            reasons=["destructive SQL statement detected"]
        )
    
    # Threshold-based action determination
    hard_block = (risk >= 60 or 
                  "graph_destructive" in categories or 
                  "system_abuse" in categories)
    
    return IntentResult(
        allowed=not hard_block,
        action="block" if hard_block else "monitor",
        risk_score=risk,
        categories=categories,
        reasons=reasons
    )
\end{lstlisting}

\subsubsection{Intent Filter Modes}

\begin{table}[ht]
\centering
\caption{Intent Filter Operating Modes}
\label{tab:intent-filter-modes}
\begin{tabular}{lp{9cm}}
\hline
\textbf{Mode} & \textbf{Behavior} \\
\hline
\texttt{enforce} & Block requests exceeding risk threshold (HTTP 400) \\
\texttt{monitor} & Annotate and log risks but never block \\
\hline
\end{tabular}
\end{table}

\subsubsection{Enforcement API}

\begin{lstlisting}[language=Python,caption={Intent filter enforcement}]
def enforce_intent(
    text: str,
    *,
    resource: str = None,
    user: Any = None,
    raise_on_block: bool = True,
) -> IntentResult:
    """Analyze and optionally enforce blocking."""
    res = analyze_intent(text)
    
    # Emit audit record
    audit_policy_decision(
        policy="intent_filter",
        subject=principal,
        action="submit",
        resource=resource,
        allowed=res.allowed,
        reason="; ".join(res.reasons),
        attributes={
            "risk_score": res.risk_score,
            "categories": res.categories
        }
    )
    
    if res.action == "block" and raise_on_block:
        raise HTTPException(
            status_code=status.HTTP_400_BAD_REQUEST,
            detail={
                "message": "Request blocked by intent filter",
                "risk_score": res.risk_score,
                "categories": res.categories,
                "reasons": res.reasons
            }
        )
    
    return res
\end{lstlisting}

\subsection{Defense in Depth}
\label{subsec:defense-in-depth}

The three data protection components work together to provide layered security:

\begin{enumerate}
    \item \textbf{Intent Filter} (Input Stage): Blocks obviously malicious requests before processing.
    \item \textbf{Output Guard} (Processing Stage): Validates and sanitizes LLM-generated queries.
    \item \textbf{PII Scrubber} (Output Stage): Removes sensitive data from responses and logs.
\end{enumerate}

This defense-in-depth approach ensures that even if one layer fails, subsequent layers provide additional protection against data leakage and unauthorized operations.

